\documentclass[10pt,a4paper]{amsart}
\usepackage[margin=19mm]{geometry}
\usepackage{amsmath,amssymb,mathtools}
\usepackage[scr=rsfs,cal=euler]{mathalfa}
\usepackage{xfrac}
\usepackage[unicode,hidelinks,bookmarks=false]{hyperref}
\newcommand{\ie}{i.e.\ }
\newcommand{\cf}{cf.\ }
\newcommand{\resp}{respectively\ }
\newcommand{\Subsection}{\subsection}
\providecommand{\nobreakdash}{\nobreak\hbox{-}}
\setlength{\emergencystretch}{2em}
\multlinegap0pt
\usepackage{bm}
\usepackage{bbm}
\usepackage{dsfont}
\usepackage{xspace}
\usepackage{cancel}
\usepackage{mathtools}
\usepackage{amsthm}
\makeatletter
\@ifundefined{lem}{\newtheorem{lem}{Lemma}}{}
\@ifundefined{thm}{\newtheorem{thm}[lem]{Theorem}}{}
\makeatother
\title[A Transmission Framework for Stark Operators with $\delta$ I]{A Transmission Framework for Stark Operators with $\delta$ Interactions on a Compact Lipschitz Interface}
\author{Masahiro Kaminaga}
\date{Source: arXiv:2509.01225v4 (CC BY 4.0)}
\begin{document}
\maketitle

\noindent\emph{Source and licence.} A Transmission Framework for Stark Operators with $\delta$ Interactions on a Compact Lipschitz Interface. Version arXiv:2509.01225v4 (CC BY 4.0), distributed under the Creative Commons Attribution 4.0 International licence, as recorded in the arXiv licence metadata for this exact version and shown on the abstract page https://arxiv.org/abs/2509.01225v4. Full rights notice: licenses/arxiv-2509.01225-CC-BY-4.0.txt.\par\smallskip

\noindent\textit{Editorial scope.} Counted unit: the Thm whose source label is \texttt{thm:resolvent-new} in the article (source0.tex lines 1093--1104, source label \texttt{thm:resolvent-new}), delivered together with the article's own complete original proof of it (source0.tex lines 1106--1238, one \texttt{proof} environment of 133 lines). No mathematics is written, repaired or re-derived: statement and proof are the article's own text line for line. Required context retained uncounted: lem:sobolev-embeddings (lines 254--273); lem:tau-dense (lines 400--412); lem:free-in-transmission (lines 620--630); lem:symmetric-HFa (lines 817--819); lem:single--layer (lines 909--940); eq:B-alpha (lines 1063--1066). Every definition, display and label of those lines is retained unchanged. Omitted from this package and not redistributed: 1--253, 274--399, 413--619, 631--816, 820--908, 941--1062, 1067--1092, 1105--1105, 1239--1575. Nothing inside a retained excerpt is omitted or altered: every retained body is delivered exactly as the article's lines are written, so no omission receipt of any kind accompanies this package. Citations: the retained excerpts carry no citation command at all, so no bibliography entry of the article is transcribed and no reference is redistributed. Reference adaptation: none was needed and none was made; every reference inside the delivered unit resolves inside it, so no unresolved reference or citation is produced. Printed slips kept literal: no printed word, symbol, hypothesis or number of the counted statement or of its proof was altered. Layout and class: the article's own class is \texttt{article} with packages including amsmath, amsfonts, amssymb, amsthm, eucal, fontenc, inputenc; the wrapper is the lane's pinned \texttt{amsart} editorial wrapper at 10pt on A4, which restates the theorem-like environments the retained text needs and adds the article's own macro definitions that the retained text needs (none), with extra wrapper packages (bm, bbm, dsfont, xspace, cancel, mathtools). The delivered numbering is the unit's own. Assets: no retained span contains a figure environment, a graphics-inclusion command, a PDF-page inclusion, a table or a TikZ picture; no image file is copied into this package, the package declares no asset dependency and no image file is redistributed with it. Licence: version arXiv:2509.01225v4 of this article is distributed under the Creative Commons Attribution 4.0 International licence, as recorded in the arXiv licence metadata for this exact version and shown on the abstract page https://arxiv.org/abs/2509.01225v4. A full rights notice is delivered at licenses/arxiv-2509.01225-CC-BY-4.0.txt. Characters: every retained excerpt is pure ASCII, so the delivered engine, XeLaTeX with its default Latin Modern text font, reports no missing character.
\begin{lem}\label{lem:sobolev-embeddings}
Let $\Sigma\subset\mathbb R^d$ be a compact Lipschitz hypersurface.
Then the embeddings
$$
H^{1/2}(\Sigma)\hookrightarrow L^2(\Sigma),
\qquad
L^2(\Sigma)\hookrightarrow H^{-1/2}(\Sigma)
$$
are continuous.
Moreover, the embedding
$$
H^{1/2}(\Sigma)\hookrightarrow L^2(\Sigma)
$$
is compact.
Consequently,
$$
H^{1/2}(\Sigma)\hookrightarrow H^{-1/2}(\Sigma)
$$
is compact.
\end{lem}

\begin{lem}\label{lem:tau-dense}
The range of the restricted trace map
$$
\tau\upharpoonright_{D_F}:D_F\to H^{1/2}(\Sigma)
$$
is dense in $H^{1/2}(\Sigma)$.
Consequently, the graph--dual adjoint
$$
\tau_D^*:H^{-1/2}(\Sigma)\to D_F'
$$
is injective.
The global adjoint $\tau_1^*:H^{-1/2}(\Sigma)\to H^{-1}(\mathbb R^d)$ is also injective.
\end{lem}

\begin{lem}\label{lem:free-in-transmission}
If $u\in D_F$, then $u\in\mathcal D_\Sigma$ and
$$
\ell_F^\Sigma u=H_{F,0}u,
\qquad
[u]=0,
\qquad
[\partial_\nu u]=0.
$$
In particular, the common trace satisfies $\gamma u=\tau u$.
\end{lem}

\begin{lem}\label{lem:symmetric-HFa}
The operator $H_{F,\alpha}$ is densely defined in $L^2(\mathbb R^d)$ and is symmetric.
\end{lem}

\begin{lem}\label{lem:single--layer}
Let $z\in\mathbb C\setminus\mathbb R$ and let $\varphi\in H^{-1/2}(\Sigma)$.
Define
$$
w:=G_F(z)\varphi=R_{F,0}(z)\tau^*\varphi.
$$
Then the following statements hold.
\begin{itemize}
\item[(i)]
$(\ell_F-z)w=\tau_1^*\varphi$ in $\mathcal D'(\mathbb R^d)$, and $(\ell_F-z)w=0$ in $\mathcal D'(\mathbb R^d\setminus\Sigma)$.

\item[(ii)]
$w\in\mathcal D_\Sigma$ and $w\in H^2_{\mathrm{loc}}(\mathbb R^d\setminus\Sigma)$.

\item[(iii)] The traces satisfy
$$
[w]=0,\qquad [\partial_\nu w]=-\varphi,
$$
and the common trace $\gamma w$ belongs to $H^{1/2}(\Sigma)$.

\item[(iv)] The map
$$
H^{-1/2}(\Sigma)\ni\varphi
\mapsto\gamma w\in H^{1/2}(\Sigma)
$$
is linear and bounded.
We denote this operator by
$$
M_F(z)\varphi:=\gamma w.
$$
\end{itemize}
\end{lem}

\begin{equation}\label{eq:B-alpha}
B_\alpha(z):=I+\alpha M_F(z)
\quad\text{on }H^{-1/2}(\Sigma).
\end{equation}

\begin{thm}\label{thm:resolvent-new}
Let $\alpha\in L^\infty(\Sigma;\mathbb R)$ and $z\in\mathbb C\setminus\mathbb R$.
Then the operator $B_\alpha(z)$ in \eqref{eq:B-alpha} is boundedly invertible on $H^{-1/2}(\Sigma)$, and
\begin{equation}\label{eq:resolvent-new}
(H_{F,\alpha}-z)^{-1}
=
R_{F,0}(z)
-
R_{F,0}(z)\tau^*(I+\alpha M_F(z))^{-1}
\alpha\tau R_{F,0}(z).
\end{equation}
\end{thm}

\begin{proof}
We first show that
$$
B_\alpha(z)=I+\alpha M_F(z):H^{-1/2}(\Sigma)
\to H^{-1/2}(\Sigma)
$$
is invertible.
The operator
$$
M_F(z):H^{-1/2}(\Sigma)\to H^{1/2}(\Sigma)
$$
is bounded.
By Lemma~\ref{lem:sobolev-embeddings}, multiplication by $\alpha$ defines a compact map
$$
\alpha:H^{1/2}(\Sigma)\to H^{-1/2}(\Sigma).
$$
Therefore $\alpha M_F(z)$ is compact on $H^{-1/2}(\Sigma)$ and $B_\alpha(z)$ is Fredholm of index zero.

It remains to prove injectivity.
Assume that
$$
B_\alpha(z)\varphi=0
$$
for some $\varphi\in H^{-1/2}(\Sigma)$, and define
$$
w:=R_{F,0}(z)\tau^*\varphi.
$$
By Lemma~\ref{lem:single--layer},
$$
(\ell_F-z)w=0
\quad\text{in }\mathbb R^d\setminus\Sigma,
\qquad
[w]=0,
\qquad
[\partial_\nu w]=-\varphi,
\qquad
\gamma w=M_F(z)\varphi.
$$
The boundary equation implies
$$
-\varphi=\alpha M_F(z)\varphi=\alpha\gamma w.
$$
Hence $w\in D(H_{F,\alpha})$ and
$$
(H_{F,\alpha}-z)w=0.
$$
By Lemma~\ref{lem:symmetric-HFa}, $(H_{F,\alpha}w,w)$ is real.
Taking imaginary parts in
$$
0=((H_{F,\alpha}-z)w,w)
$$
gives
$$
-\operatorname{Im}z\,\|w\|^2=0.
$$
Since $\operatorname{Im}z\ne0$, it follows that $w=0$.
On the other hand,
$$
(\ell_F-z)w=\tau_1^*\varphi
$$
in the sense of distributions.
Thus $\tau_1^*\varphi=0$, and Lemma~\ref{lem:tau-dense} yields $\varphi=0$.
Since $B_\alpha(z)$ is Fredholm of index zero, injectivity implies surjectivity.
Its inverse is bounded by the bounded inverse theorem.

Next let $f\in L^2(\mathbb R^d)$ and put
$$
u_0:=R_{F,0}(z)f\in D_F.
$$
By Lemma~\ref{lem:free-in-transmission}, $u_0\in\mathcal D_\Sigma$, its two jumps vanish, and $\gamma u_0=\tau u_0\in H^{1/2}(\Sigma)$.
Define
$$
\psi:=B_\alpha(z)^{-1}\alpha\tau u_0
\in H^{-1/2}(\Sigma)
$$
and set
$$
u:=u_0-G_F(z)\psi.
$$
The class $\mathcal D_\Sigma$ is linear, and both terms in this difference belong to it by Lemma~\ref{lem:free-in-transmission} and Lemma~\ref{lem:single--layer}.
Moreover, Lemma~\ref{lem:single--layer} gives
$$
[u]=0,
\qquad
[\partial_\nu u]=\psi,
\qquad
\gamma u=\tau u_0-M_F(z)\psi.
$$
The equation $B_\alpha(z)\psi=\alpha\tau u_0$ is equivalent to
$$
\psi
=\alpha\bigl(\tau u_0-M_F(z)\psi\bigr)
=\alpha\gamma u.
$$
Moreover,
$$
\ell_F^\Sigma u
=
\ell_Fu_0-zG_F(z)\psi
=
f+zu_0-zG_F(z)\psi
=
f+zu
\in L^2(\mathbb R^d).
$$
Thus $u\in D(H_{F,\alpha})$ and
$$
(H_{F,\alpha}-z)u=f.
$$
Therefore $H_{F,\alpha}-z$ is surjective.
This construction also displays the signs in the claimed resolvent formula without introducing a negative boundary density.

To prove injectivity, let $u\in D(H_{F,\alpha})$ satisfy $(H_{F,\alpha}-z)u=0$.
Since $H_{F,\alpha}$ is symmetric,
$$
0=((H_{F,\alpha}-z)u,u)
=(H_{F,\alpha}u,u)-z\|u\|^2.
$$
Taking imaginary parts gives $-\operatorname{Im}z\,\|u\|^2=0$, and hence $u=0$.
Thus $H_{F,\alpha}-z$ is bijective.

For the unique solution constructed above,
\begin{eqnarray*}
u
&=&R_{F,0}(z)f
-R_{F,0}(z)\tau^*(I+\alpha M_F(z))^{-1}
\alpha\tau R_{F,0}(z)f.
\end{eqnarray*}
Every factor on the right--hand side of \eqref{eq:resolvent-new} is bounded between the indicated spaces.
Hence the displayed solution operator is bounded on $L^2(\mathbb R^d)$.
The surjectivity and injectivity proved above show that it is the two--sided inverse of $H_{F,\alpha}-z$.
This proves \eqref{eq:resolvent-new}.
\end{proof}

\end{document}
